\chapter{Method\sv{Metod}}
\label{chap:method}
\begin{instructions}
\begin{notplan}
In this chapter, you describe how you proceeded to develop your artifact and how you designed it. The method chapter must be written so that the artifact or your results can be \textbf{reproduced}. That means that by following your instructions, someone else can end up with the same artifact or results.
\end{notplan}

In this chapter you should organize your text into the following sections:
\end{instructions}

\section{Development Process\sv{Utvecklingsprocess}}
\label{sec:devprocess}
\begin{instructions}
\textbf{Briefly} tell the reader how you proceeded in order to develop your artifact. (\emph{Please observe:} An artifact does not need to be a physical thing, it can very well be an algorithm, a mathematical theory or similar.) Examples of this can be:
\begin{itemize}
    \item a prototype
    \item a simulation
    \item a product development process
    \item a software development process,
    \item a design process (user-centered design or other)
    \item other (talk about this with your supervisor)
\end{itemize}
In any case you need to briefly explain your choice and show how you make use of your choice. You might need to mention tools for process management (e.g. git for collaboration, version control and documentation) or approaches (e.g. extreme programming for agile software development) but unless it's absolutely necessary, these only need to be mentioned very shortly and, if necessary, referenced.

You need to include how testing and assessing your artifact fit in the process. This may include
\begin{itemize}
    \item using benchmarks to test performance or
    \item generating data to test the prototype (or product) or
    \item testing usability of the prototype or
    \item a combination of these or
    \item other (talk about this with your supervisor).
\end{itemize}
\end{instructions}

\section{Technical Design\sv{Konstruktion}}
\label{sec:techdesign}
\begin{instructions}
Tell the reader \textbf{in detail} about the overall structure of your solution. It should let the reader understand how your artifact is organized, what parts it contains, how they relate to each other. This is usually called the architecture of your artifact. Use technical language and diagrams. You have to show that you are familiar with technical terminology.

It is important that you \textbf{motivate} your choice of the particular architecture (in other words, explain the technical reasons for \textbf{why} you chose this particular architecture).
\end{instructions}
